\section*{Hoofdstuk \ref{ch:b1:heat-budget} --- De warmtebalans van de aarde}

\begin{solution}{exo:b1:heat-budget:1}
$q = k\,\mathrm dT/\mathrm dz = 2.5\times\qty{0.022}{K.m^{-1}} =
\qty{0.055}{W.m^{-2}} = \qty{55}{mW.m^{-2}}$. Door \qty{1}{km^2} $=
\qty{e6}{m^2}$: \qty{5.5e4}{W}, dat is \qty{55}{kW}.
\end{solution}

\begin{solution}{exo:b1:heat-budget:2}
Continenten $\num{2.073e14}\times0.0709 = \qty{1.470e13}{W} = \qty{14.70}{TW}$;
oceanen $\num{3.028e14}\times0.1054 = \qty{31.92}{TW}$; totaal \qty{46.61}{TW}, en
de gemiddelde \omterm{def:b1:heat-budget:heat-flow}{warmtestroom} $\num{4.661e13}/\num{5.101e14} = \qty{91.4}{mW.m^{-2}}$ (de
gepubliceerde 46.7 en 91.6 verschillen alleen door de afronding van de twee gemiddelden). Het
aandeel van de zonne-instraling is $46.6/173000 = \num{2.7e-4}$, of
\qty{0.027}{\percent}.
\end{solution}

\begin{solution}{exo:b1:heat-budget:3}
Uranium: $0.99274\times\num{9.513e-5} + 0.00720\times\num{5.685e-4} =
\num{9.444e-5} + \num{4.09e-6} = \qty{9.85e-5}{W.kg^{-1}}$; het zeldzame
\ce{^{235}U} levert \qty{4}{\percent} ervan. Kalium: de massafractie van
\ce{^{40}K} is $\num{1.17e-4}\times39.96/39.10 = \num{1.196e-4}$, dus
$\num{1.196e-4}\times\num{2.847e-5} = \qty{3.40e-9}{W.kg^{-1}}$.
\end{solution}

\begin{solution}{exo:b1:heat-budget:4}
Omgerekend: $\qty{1}{\micro cal.cm^{-2}.s^{-1}} = \num{4.184e-6}\,\unit{J.s^{-1}}/\qty{e-4}{m^2}
= \qty{41.84}{mW.m^{-2}}$: oostelijke Verenigde Staten \qty{33.5}{mW.m^{-2}}, Basin and
Range-provincie \qty{58.6}{mW.m^{-2}}, Sierra Nevada \qty{16.7}{mW.m^{-2}}. Met $A_0D$ in
\unit{mW.m^{-2}}, wanneer $A_0$ in \unit{\micro W.m^{-3}} en $D$ in \unit{km} staat:
$33.5 + 3.0\times7.5 = \qty{56.0}{mW.m^{-2}}$, $58.6 + 3.0\times9.4 =
\qty{86.8}{mW.m^{-2}}$, $16.7 + 3.0\times10 = \qty{46.7}{mW.m^{-2}}$. Hetzelfde
graniet geeft in de Basin and Range een bijna twee keer zo hoge \omterm{def:b1:heat-budget:heat-flow}{warmtestroom} als in de Sierra Nevada: het verschil zit in de \omterm{def:b1:layered-earth:layers}{mantel} eronder.
\end{solution}

\begin{solution}{exo:b1:heat-budget:5}
$T(10\,\text{km}) = 10 + \frac{0.060\times10^4}{2.5} - \frac{\num{2e-6}\times10^8}{2\times2.5}
= 10 + 240 - 40 = \qty{210}{\degreeCelsius}$. Onder, $q_m = 60 - 2.0\times10 =
\qty{40}{mW.m^{-2}}$, een gradiënt van \qty{16}{K.km^{-1}}, en
$T(35\,\text{km}) = 210 + 16\times25 = \qty{610}{\degreeCelsius}$.
\end{solution}

\begin{solution}{exo:b1:heat-budget:6}
Schalie: gradiënt $12.0/0.300 = \qty{40}{K.km^{-1}}$, $q = 1.5\times0.040 =
\qty{60}{mW.m^{-2}}$. Graniet: $10.0/0.500 = \qty{20}{K.km^{-1}}$,
$q = 3.0\times0.020 = \qty{60}{mW.m^{-2}}$. $R(300) = 300/1.5 =
\qty{200}{m^2.K.W^{-1}}$, $R(800) = 200 + 500/3.0 = \qty{366.7}{m^2.K.W^{-1}}$.
De drie punten $(0, 12.0)$, $(200, 24.0)$, $(366.7, 34.0)$ liggen op één lijn, met
helling $22.0/366.7 = \qty{0.0600}{W.m^{-2}}$: de \omterm{def:b1:heat-budget:heat-flow}{warmtestroom} is in beide gesteenten dezelfde,
\qty{60}{mW.m^{-2}}; alleen de gradiënt verandert, in omgekeerde
verhouding tot de geleidbaarheid.
\end{solution}

\begin{solution}{exo:b1:heat-budget:7}
$\tau = L^2/\kappa$: \qty{1}{m}, \qty{e6}{s} of \num{12}~dagen; \qty{10}{m},
\qty{e8}{s} of \num{3.2}~jaar; \qty{100}{km}, \qty{e16}{s} of
\qty{317}{Myr}; \qty{2900}{km}, \qty{8.4e18}{s} of \qty{267}{Gyr}. Op
\qty{10}{m} is de dagelijkse cyclus helemaal niet voelbaar, en de jaarlijkse, die korter duurt dan de drie jaar die nodig zijn om \qty{10}{m} te doorkruisen, slechts zeer zwak: alleen veranderingen die jaren of langer duren komen door. De laatste tijd is $267/4.54 = 59$ keer de leeftijd van de aarde.
\end{solution}

\begin{solution}{exo:b1:heat-budget:8}
(a) $1646 - 0.54\times50 = \qty{1619}{K}$, dat is \qty{1346}{\degreeCelsius},
de gepubliceerde \qty{1350}{\degreeCelsius} afgerond. (b) Gemiddelde gradiënt
$(0.54 + 0.36)/2 = \qty{0.45}{K.km^{-1}}$, dus $1646 + 0.45\times360 =
\qty{1808}{K}$, \qty{9}{K} (\qty{0.5}{\percent}) boven \qty{1799}{K}: de gradiënt neemt niet lineair af, maar de schatting ligt dicht erbij.
\end{solution}

\begin{solution}{exo:b1:heat-budget:9}
(a) $G^2 = T_0^2/(\pi\kappa t)$ geeft $t = T_0^2/(\pi\kappa G^2)$.
(b) $G = 0.02\times\tfrac59\,\unit{K}/\qty{0.3048}{m} = \qty{0.03645}{K.m^{-1}} =
\qty{36.45}{K.km^{-1}}$; $t = 3900^2/(\pi\times\num{1.2e-6}\times0.03645^2) =
\num{1.521e7}/\num{5.009e-9} = \qty{3.04e15}{s}$, dat is \qty{96}{Myr},
Kelvins getal. (c) $t = 1300^2/(\pi\times\num{e-6}\times0.036^2) =
\qty{4.15e14}{s} = \qty{13}{Myr}$: een realistischer begintemperatuur maakt de geleidende aarde nog jonger.
\end{solution}

\begin{solution}{exo:b1:heat-budget:10}
(a) Na de erosie is het nieuwe oppervlak de oude diepte $e$, waar de warmteproductie $A_0' = A_0e^{-e/D}$ was en de \omterm{def:b1:heat-budget:heat-flow}{warmtestroom} $q(e) = q_m +
A_0De^{-e/D} = q_m + A_0'D$; gemeten vanaf het nieuwe oppervlak neemt de productie nog steeds af als $A_0'e^{-z'/D}$. De pluton beweegt langs dezelfde lijn omlaag.
(b) De laag is nu $h - e$ dik met dezelfde $A$, dus $q_0' = q_m + A(h - e)$:
op een lijn met helling $h - e$ door $q_m$. (c) Alleen het exponentiële model: in het uniforme-lagenmodel zouden plutons die tot verschillende diepten zijn geërodeerd op lijnen met verschillende hellingen liggen, en de punten van één provincie zouden verspreid zijn. Dat er één lijn wordt waargenomen, wijst op een productie die ongeveer exponentieel met de diepte afneemt.
\end{solution}

\begin{solution}{exo:b1:heat-budget:11}
(a) $\tau_P = RL/(3\kappa)$: voor \qty{50}{km}, $\num{6.4e6}\times\num{5e4}/\num{3e-6}
= \qty{1.07e17}{s} = \qty{3.4}{Gyr}$; voor \qty{100}{km}, \qty{6.8}{Gyr}.
(b) $q = kT_i/L = 3\times1300/\num{5e4} = \qty{78}{mW.m^{-2}}$ bij het begin,
en $78\,e^{-2/3.38} = \qty{43}{mW.m^{-2}}$ na \qty{2}{Gyr}. (c)
$3\times1300/\sqrt{\pi\times\num{e-6}\times\num{6.31e16}} = 3900/\num{4.45e5} =
\qty{8.8}{mW.m^{-2}}$: de goed doorgeroerde aarde verliest na twee miljard jaar nog vijf keer meer warmte. Haar hele binnenste voedt het oppervlak, terwijl in een geleidende aarde slechts een buitenste schil dat doet.
\end{solution}

\begin{solution}{exo:b1:heat-budget:12}
(a) $\sqrt{\pi\kappa\times\qty{1}{Myr}} = \sqrt{\pi\times\num{e-6}\times\num{3.156e13}}
= \qty{9957}{m}$, dus $C = 3\times1336/9957 = \qty{0.403}{W.m^{-2}}$ voor
$t$ in Myr, \qty{403}{mW.m^{-2}.Myr^{1/2}}: \qty{82}{\percent} van de op de gegevens gefitte
490. Het model heeft een grotere $k\,\Delta T$ nodig (ongeveer \qty{4900}{W.m^{-1}}): een hogere geleidbaarheid van de \omterm{def:b1:layered-earth:layers}{mantel}, of een heter binnenste, zoals langs de adiabaat.
(b) $490/\sqrt1 = 490$, $490/\sqrt{10} = 155$ en $490/\sqrt{100} =
\qty{49}{mW.m^{-2}}$. (c) $z = 2.326\sqrt{\num{e-6}\times\num{3.156e15}} =
2.326\times\qty{56180}{m} = \qty{131}{km}$.
\end{solution}

\begin{solution}{pb:b1:heat-budget:1}
\textbf{1.} Gradiënt $= (T_{500} - T_{100})/\qty{0.4}{km}$: A \qty{15.25}{K.km^{-1}},
B \num{19.50}, C \num{19.50}, D \num{25.50}, E \qty{28.25}{K.km^{-1}}.

\textbf{2.} Het boren verstoort de temperatuur van het gesteente (de boorvloeistof koelt of verwarmt de wanden), en de storing duurt weken tot maanden om weg te ebben. De bovenste tientallen meters volgen de seizoenscyclus en recente veranderingen van de oppervlaktetemperatuur.

\textbf{3.} $q_0 = k\times$gradiënt: A $2.9\times15.25 = \qty{44.225}{mW.m^{-2}}$,
B \qty{50.70}{mW.m^{-2}}, C \qty{60.45}{mW.m^{-2}}, D \qty{68.85}{mW.m^{-2}},
E \qty{84.75}{mW.m^{-2}}.

\textbf{4.} A tot en met D liggen onder het continentale gemiddelde, E erboven, met
\qty{13.9}{mW.m^{-2}}, of \qty{20}{\percent}.

\textbf{5.} De \omterm{def:b1:heat-budget:heat-flow}{warmtestroom} is geleidbaarheid maal gradiënt: C geleidt beter (3.1
tegen 2.6), dus dezelfde gradiënt voert \qty{19}{\percent} meer warmte af.

\textbf{6.} De isothermen volgen het reliëf maar gedeeltelijk: onder een top stijgen zij minder dan het oppervlak, dus liggen zij daaronder verder uit elkaar, en wordt de gemeten gradiënt verlaagd (onder een dal verhoogd). Een correctie voor het reliëf is nodig voordat de plutons met elkaar worden vergeleken.

\textbf{7.} $\rho C_U H_U$: A $2650\times\num{2.0e-6}\times\num{9.85e-5} =
\qty{0.522}{\micro W.m^{-3}}$, B 0.914, C 1.305, D 1.827, E 2.610.

\textbf{8.} Thorium: A 0.558, B 0.976, C 1.255, D 1.742, E 2.439.

\textbf{9.} Kalium: A $2650\times0.030\times\num{3.40e-9} = 0.270$, B 0.315,
C 0.360, D 0.378, E 0.405.

\textbf{10.} $A_0$ = 1.35, 2.20, 2.92, 3.95 en \qty{5.46}{\micro W.m^{-3}}.
Kalium draagt het minst bij, van \qty{20}{\percent} van het totaal in A tot
\qty{7}{\percent} in E; uranium en thorium delen de rest vrijwel gelijk.

\textbf{11.} Ja: de vijf punten stijgen gestaag, binnen ongeveer
\qty{1.5}{mW.m^{-2}} van een rechte lijn.

\textbf{12.} $\bar A_0 = 15.88/5 = 3.176$, $\bar q_0 = 308.975/5 = 61.795$;
$\sum(A_0 - \bar A_0)^2 = 10.1681$; $\sum(A_0 - \bar A_0)(q_0 - \bar q_0) =
101.146$.

\textbf{13.} $D = 101.146/10.1681 = 9.95$ (in \unit{mW.m^{-2}} per
\unit{\micro W.m^{-3}}, dat is in km): $D = \qty{9.95}{km}$, en $q_m = 61.795 -
9.95\times3.176 = \qty{30.20}{mW.m^{-2}}$.

\textbf{14.} De helling, ongeveer \qty{10}{km}, is die van de drie provincies; de gereduceerde \omterm{def:b1:heat-budget:heat-flow}{warmtestroom}, \qty{30}{mW.m^{-2}}, ligt dicht bij die van de oostelijke Verenigde Staten, \qty{33}{mW.m^{-2}}: een rustige provincie met een koele \omterm{def:b1:layered-earth:layers}{mantel} eronder, anders dan de Basin and Range-provincie.

\textbf{15.} $q_m/k = \qty{12}{K.km^{-1}}$ en $A_0D^2/k = \num{5.46e-6}\times\num{e8}/2.5
= \qty{218.4}{K}$: $T(z) = 10 + 12z + 218.4\,(1 - e^{-z/10})$, $z$ in km, $T$ in
\unit{\degreeCelsius}.

\textbf{16.} $q(z) = 30 + 54.6\,e^{-z/10}$: \qty{50.1}{mW.m^{-2}} op \qty{10}{km},
\num{37.4} op \qty{20}{km} en \qty{31.7}{mW.m^{-2}} op \qty{35}{km}.

\textbf{17.} $30/84.6 = \qty{35}{\percent}$: bijna twee derde van de \omterm{def:b1:heat-budget:heat-flow}{warmtestroom}
boven E wordt geproduceerd in haar radiogene laag.

\textbf{18.} $T(10) = 10 + 120 + 218.4\times0.632 = \qty{268}{\degreeCelsius}$.

\textbf{19.} $T(20) = 10 + 240 + 218.4\times0.865 = \qty{439}{\degreeCelsius}$.

\textbf{20.} Onder A is $A_0D^2/k = 1.35\times100/2.5 = \qty{54.0}{K}$:
$T(35) = 10 + 420 + 54.0\times0.970 = \qty{482}{\degreeCelsius}$.

\textbf{21.} Door $A_0D^2/k$: \qty{54}{K} onder A en \qty{218}{K} onder E. De diepe \omterm{def:b1:heat-budget:heat-flow}{warmtestroom} is dezelfde; E's \omterm{def:b1:layered-earth:layers}{korst} is vier keer rijker aan warmtebronnen, en de warmte die zij produceert moet door hetzelfde gesteente naar boven worden geleid, wat over de hele weg een steilere gradiënt vergt.

\textbf{22.} $1300 = 10 + 218.4 + 12z$ geeft $z = \qty{89.3}{km}$ onder E;
$1300 = 10 + 54.0 + 12z$ geeft \qty{103.0}{km} onder A.

\textbf{23.} Beide liggen binnen het bereik van de continentale \omterm{def:b1:layered-earth:lithosphere}{lithosfeer}, van honderd kilometer of minder onder jonge gebieden tot ongeveer \qty{250}{km} onder cratons, aan de dunne kant daarvan: bij een gereduceerde \omterm{def:b1:heat-budget:heat-flow}{warmtestroom} van \qty{30}{mW.m^{-2}} en deze geleidbaarheid bereikt de conductieve \omterm{def:b1:heat-budget:geotherm}{geotherm} manteltemperaturen op ongeveer \qty{100}{km}; onder de sterker radioactieve pluton \qty{14}{km} eerder.

\textbf{24.} $T(35) = 10 + 420 + 218.4\times(1 - e^{-3.5}) = 430 + 211.8$: de
\omterm{def:b1:layered-earth:discontinuity}{Moho} onder pluton E ligt bij
\textbf{\qty[text-series-to-math=true]{642}{\degreeCelsius}}.
\end{solution}
